\begin{textsample}{POS Dim 2 – mg – Score 24.00 – mg/mg\_inf\_e\_ef\_intro\_1.txt}  \label{ex:f2_pos_199}
1.1 Apresentação Caro(a) Educador(a), É com grande satisfação que lhe apresentamos o \textbf{Currículo} Referência de Minas Gerais, documento elaborado a partir dos fundamentos educacionais expostos na nossa Constituição Federal ( CF/1988 ), na Lei de \textbf{Diretrizes} e Bases da Educação Nacional ( LDB 9394/96 ), no Plano Nacional de Educação ( PNE/2014 ), na Base Nacional Comum Curricular ( BNCC/2017 ), de acordo com as determinações e orientações do Conselho Nacional de Educação expressos por meio de pareceres/resoluções e a partir do reconhecimento e da valorização dos diferentes povos, culturas, crenças, territórios e tradições existentes em nosso Estado.
Convidamos você para trilhar pelo histórico da construção do documento, pela concepção de \textbf{currículo} e sua prática em sala de aula até chegar aos eixos estruturadores que compõem o mesmo.
É importante que você se sinta parte deste texto, uma vez que é você quem o colocará em prática na sala de aula.
Destacamos que o presente documento é resultado do \textbf{regime} de colaboração estabelecido entre a Secretaria de Estado de Educação de Minas Gerais – SEE/MG e a União Nacional dos \textbf{Dirigentes} \textbf{Municipais} de Educação de Minas Gerais, seccional Minas Gerais - UNDIME/MG.
Neste processo, foram considerados e estudados os documentos curriculares já presentes em diferentes \textbf{redes} ( \textbf{estadual} e \textbf{municipais} ) como fonte de inspiração para a elaboração de um \textbf{currículo} que possa ser referência em todo o Estado.
A construção, portanto, do \textbf{Currículo} Referência de Minas Gerais de forma dialogada entre o Estado e os Municípios é uma oportunidade singular para o \textbf{fortalecimento} do \textbf{regime} de colaboração – previsto desde a Constituição Federal ( 1988 ) – e para avançarmos na consolidação de um Sistema Integrado de Educação Pública - SIEP, como vem sendo debatido em nosso Estado.
A defesa de um sistema de educação único fundamenta -se no reconhecimento desse sistema como uma das formas de manter a \textbf{integralidade} do atendimento e a \textbf{oferta} de uma educação pública inclusiva, com \textbf{qualidade} e \textbf{equidade}.
Em outras palavras, trata -se de não distinguir a qual \textbf{rede} um estudante se vincula ao longo da \textbf{trajetória} escolar.
O que se deve garantir é a \textbf{oferta} de um ensino de \textbf{qualidade} e de oportunidades de formação e transformação social diversificadas, que zelem pelo direito à aprendizagem – como destacaremos ao longo de todo o documento.
Compreendemos, assim, que é necessário “ vencer as amarras ” institucionais e culturais existentes e avançar no \textbf{fortalecimento} da democracia, na colaboração entre as \textbf{redes} e nas oportunidades aos nossos estudantes.
Ter um \textbf{currículo} referência para o Estado materializa essa proposta e potencializa as relações entre as secretarias, escolas, professores e estudantes.
Esperamos que deste documento, que ora lhe apresentamos e que orienta a elaboração dos planos e ações educacionais em Minas Gerais, surjam, da prática de cada educador e de cada escola, formas de implementação de uma educação inclusiva, igualitária e democrática.
Ao final, educador(a), você deverá reconhecer a relação deste documento com o seu dia a dia na escola, ter consciência e compreensão dos \textbf{preceitos} aqui tratados e a capacidade de trabalhá -los na sala de aula e na escola, a partir de sua autonomia, de forma dialogada com seus pares, com a gestão escolar e com os seus estudantes, fortalecendo os princípios do direito à aprendizagem de \textbf{qualidade}, da \textbf{equidade}, do reconhecimento e da valorização das diversidades, da inclusão e da gestão democrática e participativa, com vistas a promover a educação em sua \textbf{integralidade}. 1.2 O processo de construção do \textbf{Currículo} Referência de Minas Gerais Em um momento histórico de definição de uma Base Nacional Comum Curricular e elaboração de um \textbf{Currículo} Referência que \textbf{atenda} a todo o Estado, o \textbf{Regime} de Colaboração passa a ser central no cenário educacional e um desafio para Minas Gerais, em razão da sua extensão, número de \textbf{municípios} e escolas.
Mas em um esforço conjunto para reunir a imensa “ Minas Gerais ” e construir um documento coletivo, a seccional da União Nacional dos \textbf{Dirigentes} \textbf{Municipais} de Educação – UNDIME/MG, as escolas \textbf{privadas} de educação e a Secretaria de Estado de Educação – SEE/MG passam a colaborar lado a lado para redação deste documento, entendendo, fundamentalmente, que os estudantes transitam entre as \textbf{redes} ao longo da vida, ora em escolas \textbf{municipais}, ora em escolas \textbf{estaduais}, ora em escolas privadas, bem como, transitam entre os territórios, daí a importância de uma parte comum nos \textbf{currículos}.
Se o documento curricular pretende garantir de fato os direitos de aprendizagem, torna -se impossível fragmentar a vida escolar de nossos estudantes.
Portanto, pretendemos em colaboração, garantir \textbf{trajetórias} de sucesso acadêmico, somadas ao desenvolvimento integral das crianças, \textbf{adolescentes}, jovens e adultos que estão e estarão na educação pública, buscando, através de instâncias decisórias estabelecidas, o diálogo permanente entre União, Estado e Municípios.
Minas Gerais é o Estado brasileiro com maior número de \textbf{municípios} ( 853 ), representando 15 \% do total do país ( 5570 \textbf{municípios} ).
O Estado é um retrato quase sempre fiel da realidade brasileira, com 10 \% ( 20.7 milhões ) da população nacional ( 209.3 milhões ), representando a grande diversidade regional, econômica, política e social.
Em termos educacionais, nosso Estado conta com 16.151 escolas, das quais 3.622 são \textbf{estaduais}, 8.751 \textbf{municipais} e 3.778 privadas, distribuídas em 47 regionais de ensino ( SRE ), e 4.032.949 de estudantes matriculados, sendo que 86 \% deles estão na \textbf{rede} pública.
Com a maioria das escolas e das matrículas pertencentes à \textbf{rede} pública, garantir uma educação de \textbf{qualidade} com \textbf{equidade} é princípio norteador das \textbf{políticas} públicas de educação nas \textbf{redes} \textbf{municipais} e \textbf{estadual}.
Localizado na região Sudeste, o Estado de Minas Gerais é o quarto em área territorial do país com uma área de 586.520,732 km².
E com uma população estimada de 21.119.536 habitantes distribuídos em 853 \textbf{municípios}, é o segundo mais populoso do país, o que contribui significativamente para sua grande pluralidade regional.
A diversidade regional do Estado de Minas Gerais é resultado de um processo histórico de ocupação do território marcado por diferentes fatores, desde aqueles de ordem socioeconômica até os naturais de clima e vegetação.
Essa diversidade se traduz no que podemos entender como várias “ Minas Gerais ” dentro dos limites do Estado, exigindo, portanto, diferentes formas de abordagem e atuação sobre a realidade mineira.
De fato, a efetividade de qualquer iniciativa parte necessariamente da compreensão da realidade para a qual se propõe.
Para elaboração do novo \textbf{currículo} que \textbf{atenda} tamanha diversidade foi estabelecido um modelo de governança dinâmico, capaz de lidar com as particularidades de Minas Gerais e com as diversas \textbf{entidades} que atuam diretamente para melhoria da educação pública.
E para isso, foi estabelecida uma Comissão Estadual, com representações políticas de \textbf{órgãos} e \textbf{entidades}, um Comitê Executivo para condução e tomada de decisão, uma Coordenação \textbf{Técnica} para encaminhamento dos trabalhos e Grupos de Trabalho de \textbf{Currículo} para redação do documento, conforme Figura 1.

% matched lemmas: adolescente, atender, currículo, diretriz, dirigente, entidade, equidade, estadual, fortalecimento, integralidade, municipal, município, oferta, política, preceito, privar, qualidade, rede, regime, trajetória, técnico, órgão
\end{textsample}
